\documentclass[11pt]{article}
\usepackage[margin=1in]{geometry}
\usepackage[T1]{fontenc}
\usepackage{lmodern}
\usepackage{microtype}
\usepackage{parskip}

\title{David Chivers academic voice}
\date{}

\begin{document}
\maketitle

This PDF shows a short before-and-after example of the academic voice memory working on a rough paragraph, and identifies the specific voice rules applied in the rewrite.

\section*{Before}

A lot of local residents do not like new housing and this is a big reason prices are high. Older homeowners often oppose development because they want to protect the value of their houses, and younger people are then left worse off. This can have major long-run effects on affordability and inequality.

\section*{After}

Local control over zoning gives incumbent homeowners direct influence over the supply of new housing. As additional construction lowers scarcity rents and can alter neighbourhood composition, older homeowners often support tighter restrictions, while younger renters benefit from expansion. We therefore describe housing supply not as a purely technological margin, but as the outcome of a political conflict over prices and access. In this setting, demographic ageing shifts political support toward tighter regulation and raises prices by enlarging the coalition in favour of scarcity.

\section*{What from the memory file was imposed}

\begin{itemize}
\item The paragraph defines the institutional setting first: local control over zoning.
\item The mechanism is stated in causal order: control over supply, then incentives, then political support, then prices.
\item Winners and losers are named explicitly: older homeowners and younger renters.
\item The language is measured and academic. Vague phrases such as ``a lot'' and ``big reason'' are removed.
\item The paragraph ends with the economic implication rather than with a generic claim about importance.
\end{itemize}

\end{document}
